\section{Experiments}
\label{sec:experiments}

\subsection{Selecting the training regime}
\label{sec:protocol}

Grokking on modular addition occurs only within a narrow band of
hyperparameters, and the band depends on the architecture. Since the grokking
step is our principal measurement, the regime was fixed by a sweep before any
architectural comparison, rather than assumed from the literature.

Table~\ref{tab:regime} reports the dense MLP under three training fractions
and two weight-decay settings. The pattern is unambiguous. At a $30\%$
training fraction the network memorises completely --- training accuracy
reaches $1.0$ within two thousand steps --- while test accuracy creeps to
$5.5\%$ by step $30\,000$ and never generalises within budget. At $80\%$ the
transition occurs at step $2000$. Weight decay $3.0$ destroys generalisation
at every fraction.

\begin{table}[t]
\centering
\caption{Dense MLP baseline: grokking step and final test accuracy across
training regimes. A grokking step of ``---'' means the $95\%$ threshold was
never reached within $30\,000$ steps.}
\label{tab:regime}
\begin{tabular}{llrr}
\toprule
Train fraction & Weight decay & Grokking step & Final test acc. \\
\midrule
0.3 & 1.0 & --- & 0.055 \\
0.3 & 3.0 & --- & 0.000 \\
0.5 & 1.0 & 8500 & 0.985 \\
0.5 & 3.0 & --- & 0.554 \\
0.8 & 1.0 & 2000 & 0.995 \\
0.8 & 3.0 & --- & 0.775 \\
\bottomrule
\end{tabular}
\end{table}

We initially selected a $30\%$ training fraction, following the transformer
setting of \citet{nanda2023progress}. That choice turned out to be wrong for
this architecture: with one-hot inputs and no embedding layer, each pair
activates a unique pair of coordinates, so memorisation is cheap and $30\%$ of
the grid gives the network no reason to search for the generalising algorithm.
The MLP reference of \citet{bussmann2023modadd} uses $80\%$, and that is the
regime adopted here. We report the sweep rather than silently adopting the
final setting, because the failure mode it reveals --- an architecture that
looks incapable of generalising when it is merely under-supplied with data ---
recurs below for the Mixture-of-Experts layers and is easy to mistake for an
architectural result.

\subsection{Main comparison}

The principal experiment holds the task, data, split and optimiser fixed and
varies only the representation of the expert tensor. The six mixture layers are trained at $N \in \{16, 64, 256, 1024, 4096\}$
composed experts, with three seeds per configuration, and the dense MLP is
trained under the same protocol as an $N$-independent reference, giving the
grokking step, final accuracy, effective frequency count, spectral purity and
router--basis overlap for each. The
unfactorised layer is capped at $N = 1024$ for the reason given in
Section~\ref{sec:method}.

Three questions are asked of these runs. Does the number of experts help or
hurt, and does the answer depend on the factorisation? Does the representation
that a factorised layer learns differ from a dense layer's in the basis the
task defines? And where in the layer --- router or shared basis --- does the
structure that implements the algorithm reside?

\subsection{Ablations}
\label{sec:ablations}

Four ablations isolate mechanisms that would otherwise remain conjectural.

\paragraph{Router-key weight decay.} Whether the router keys are subject to
weight decay. Under one-hot inputs the keys form a lookup table of $4 n P$
free parameters addressed directly by the operand value. If the usual
convention of exempting router parameters from decay is followed, that table
constitutes an unregularised memorisation mechanism, and the prediction is
that generalisation should fail at large $N$ and be rescued by decaying the
keys.

\paragraph{Router input split.} Whether the two halves of the router input are
the two operands (the natural split), or an interleaving or random permutation
of the input coordinates. Under the natural split the factorisation of the
router matches the factorisation of the task, which arguably hands the model
the structure it is supposed to discover. The permuted variants test whether
the observed behaviour survives when that alignment is removed.

\paragraph{Auxiliary loss weight.} The coefficient $\lambda \in \{0, 10^{-3},
10^{-2}\}$ on MONET's uniformity and ambiguity terms. These were calibrated
for language modelling, where cross-entropy never approaches zero. On an
algorithmic task the cross-entropy does approach zero after memorisation, at
which point the auxiliary terms can exceed it in magnitude and begin to
dominate the gradient.

\paragraph{Training fraction.} Fractions $\{0.4, 0.6, 0.8\}$ for a
shared-basis and a product-key layer, locating the data threshold at which
each architecture stops generalising.

\subsection{What we deliberately do not measure}

One family of metrics is absent from this study by design. Both MONET and
\citet{oldfield2024mumoe} score expert specialisation over classes or domains:
an expert belongs to a domain if its routing score is skewed towards it, or if
ablating it costs accuracy on one class in particular. Neither criterion is
meaningful on single-task modular addition. There is no domain structure to be
skewed towards, and every frequency participates in predicting every value of
$c$, so an expert cannot be class-specialised even in principle. Applying
those criteria here would produce a null result that says nothing about the
criteria and nothing about the models.

Our measurements are therefore frequency-based throughout: purity, effective
frequency count, and router--basis overlap, all defined against the task's
known solution rather than against a class partition. The claim in
Section~\ref{sec:res-router-consequence} is correspondingly narrow --- that
the router stops carrying the structure that implements the computation, which
is a statement about where the mechanism lives, not a direct test of MONET's
attribution procedure.

Running that direct test needs a task with both a known algorithm and a genuine
partition of the input distribution. The mixed dataset of
Section~\ref{sec:method} was built for it: addition and multiplication examples
distinguished by a task token, where specialisation by subtask is possible in
principle and class-based criteria can be applied and checked against ablation.
Section~\ref{sec:res-mixed} reports that test. It matters in both directions:
it tells us what routing-score attribution recovers when a partition is
present, and it establishes that the frequency-based null results above are
properties of the single task rather than of our instruments.